\section{Exceptions, Transitions, and Cases That Expose the Rule}

\subsection{Married former non-Catholic ministers}

The Apostolic See has admitted some married former Anglican and other non-Catholic ministers to Latin priestly ordination after individualized discernment. The complementary norms for personal ordinariates similarly permit a case-by-case petition for a married former Anglican minister while retaining celibacy as the norm and excluding married men from episcopal ordination. These provisions demonstrate that celibacy is not intrinsic to valid priesthood. They do not create a general entitlement, abolish canon 277, or make the Eastern married presbyterate an ``exception'' to Latin law.

The candidate's wife and family are not incidental facts. Validity and stability of marriage, ecclesial reception, children's welfare, finances, formation, and the content of any rescript all require attention. A former minister's prior professional identity does not create a claim to ordination.

\subsection{Transfer between Churches \emph{sui iuris}}

A Catholic does not change canonical enrollment merely by attending another rite. A transfer of ascription follows the law and cannot be used as a device to evade the discipline of one's own Church. A married man's possible admission to orders after a lawful transfer still depends on the receiving Church's particular law, the competent authority, suitability, and all other requirements. Ritual preference is not a canonical shortcut.

\subsection{Religious clerics}

A public perpetual vow of chastity in a religious institute and the clerical obligation attached to orders are related but distinct. Canon 599 defines the evangelical counsel of chastity undertaken for the kingdom; canon 1037 recognizes perpetual religious profession when it does not require a second public assumption of celibacy before diaconate. Departure or dismissal from an institute does not erase sacred ordination, and a dispensation from vows must not be confused with a dispensation from clerical celibacy or the impediment of orders.

\subsection{Death, nullity, dissolution, and civil divorce}

These events have different legal effects. Death ends the marital bond but not the impediment arising from orders. A declaration of nullity says that a valid marriage was not established; it does not erase ordination. A privilege or dissolution may end a bond under its own conditions but does not itself authorize a cleric's new marriage. Civil divorce regulates civil status but ordinarily leaves the canonical bond intact. In every case, one must identify the act, its forum, and the separate requirements of clerical law.

\subsection{Dispensation and loss of clerical state}

Rescripts must be read, not guessed. Dismissal or rescripted loss of the clerical state, dispensation from celibacy, dispensation from the impediment of sacred orders, release from religious vows, and permission concerning a proposed marriage are distinct juridic effects even when coordinated in one Roman proceeding. Canon 291 makes the central Latin rule plain: loss of state alone does not dispense celibacy.

\begin{threestable}{Event}{What it changes}{What it does not automatically change}
Wife's death & Ends the marriage; the deacon becomes widowed. & Sacred orders, clerical state, or the impediment to a new marriage.\\
Civil divorce & Civil marital status and related civil consequences. & The canonical bond, sacred orders, or permission to marry.\\
Declaration of nullity & Judicial determination that the alleged marriage was invalid from the start. & Sacred orders or its impediment.\\
Loss of clerical state & Clerical rights and obligations as specified in law or rescript. & Valid ordination; in Latin law, celibacy unless separately dispensed.\\
Dispensation from celibacy/orders impediment & The effects exactly granted by competent authority. & A guarantee that another proposed marriage is otherwise valid or pastorally prudent.\\
Change of Church \emph{sui iuris} & Canonical ascription when lawfully completed. & Automatic eligibility for orders or evasion of formation and particular law.\\
\end{threestable}
